% from main.tex
%%%%%%%%%%  scan idx 507  |  folio 494  %%%%%%%%%%
module des différentielles \(\Omega^1_{B/\mathcal O_P}\) est isomorphe à \(B(-1)\otimes_{\mathcal O_P}\mathcal O_P(1)\) car on peut définir une \(S(E)\)-dérivation \(d\) :
\[
S(E)[T,T^{-1}]
\longrightarrow
S(E)[T,T^{-1}]\otimes_{S(E)}S(E)\cdot dT
\]
de façon à satisfaire la propriété universelle qui définit \(\Omega^1_{B/\mathcal O_P}\) en posant \(d(aT^n)=naT^{n-1}\otimes dT\) ; de plus, \(d\) est compatible aux graduations si on attribue le degré \(1\) au facteur \(S(E)\cdot dT\).

Soit \(V\) le \(Y\)-fibré vectoriel défini par \(E\). Alors \(\Omega^1_{Z/Y}\) est la restriction à \(Z\) de \(\Omega^1_{V/Y}\). Or \(V\) est le \(Y\)-schéma affine défini par \(S(E)\), donc \(\Omega^1_{V/Y}\) est défini par \(S(E)\otimes_{\mathcal O_Y}E\) (car \(E\) est localement libre de type fini), donc \(\Omega^1_{V/Y}\) est égal à \(g^*(E)\), et par suite, \(\Omega^1_{B/Y}=f^*(E)\otimes_{\mathcal O_P}B\). Pour voir sa graduation, on regarde localement, et on voit aussitôt que \(f^*(E)\) doit être considéré comme étant de degré \(1\).

En prenant les termes de degré \(0\) dans la suite exacte (5.2.1), on trouve la suite exacte annoncée.

\textbf{Remarque 5.2.1.} En utilisant les résultats généraux sur l'étude différentielle des foncteurs de Hilbert ([2], 5), on obtient directement un isomorphisme \(\Omega^1_{P/Y}\xrightarrow{\sim}F(-1)\) d'après la caractérisation fonctorielle donnée de \(g^*(\Omega^1_{P/Y})\) pour tout \(Y\)-morphisme \(g:Y'\to P\).

\textbf{5.3.} D'après 5.2, on a :
\[
 T_f=\cl(\Omega^1_{P/Y})=g^*(E)\cdot\xi^{-1}-1.
\]
Par conséquent :
\[
 \widetilde c(T_f)=\widetilde c(E)\cdot\widetilde c(\xi^{-1})-\widetilde c(1).
\]

